\providecommand{\C}{\mathbb{C}}
\providecommand{\Z}{\mathbb{Z}}
\providecommand{\Cstar}{\mathbb{C}^{\times}}
\providecommand{\cO}{\mathcal{O}}
\providecommand{\Db}{D^{b}}
\providecommand{\Heq}{H^{*}_{\mathrm{eq}}}
\providecommand{\HBM}{H^{\mathrm{BM}}}
\providecommand{\Mred}{\mathcal{M}^{\mathrm{red}}_{\DT}}
\providecommand{\DT}{\mathrm{DT}}
\providecommand{\CoHA}{\mathrm{CoHA}}
\providecommand{\Yplus}{Y^{+}}
\providecommand{\Ypluscy}{Y^{+}(K3\times E)}
\providecommand{\PhiTenUn}{\Phi_{10}^{\mathrm{un}}}
\providecommand{\PhiTenOP}{\Phi_{10}^{\mathrm{OP}}}
\providecommand{\DeltaFive}{\Delta_{5}}
\providecommand{\DFive}{D_{5}}
\providecommand{\gDeltaFive}{\mathfrak{g}_{\DeltaFive}}
\providecommand{\gBPS}{\mathfrak{g}_{\mathrm{BPS}}}
\providecommand{\phiZO}{\phi_{0,1}}
\providecommand{\kBKM}{\kappa_{\mathrm{BKM}}}
\providecommand{\kch}{\kappa_{\mathrm{ch}}}
\providecommand{\kchHeis}{\kappa_{\mathrm{ch}}^{\mathrm{Heis}}}
\providecommand{\kcat}{\kappa_{\mathrm{cat}}}
\providecommand{\kfib}{\kappa_{\mathrm{fiber}}}
\providecommand{\Mbps}{\mathcal{M}^{+}_{\mathrm{eff}}}
\providecommand{\AutS}{\mathrm{Aut}_{s}}
\providecommand{\NE}{\mathrm{NE}}
\providecommand{\Etr}{E_{\mathrm{tr}}}
\providecommand{\Tred}{T_{\mathrm{red}}}
\providecommand{\smult}{\operatorname{smult}}

\section*{The primitive $N=1$ character theorem on $K3\times E$}
\label{sec:n1-character-k3e}

Set $X=K3\times E$. The four numerical objects attached to $X$ are the
reduced Donaldson--Thomas scalar, the positive Hall half, the PBW
primitive BPS Lie algebra, and the Borcherds--Gritsenko--Nikulin
denominator algebra $\gDeltaFive$. The unconditional common statement is
the character identity. The compact Hall--Drinfeld double and the
bracket identification require separate recognition data.

\subsection*{Primitive data}

Let
\[
 \gamma=(\beta_h,d)\in H_2(K3,\Z)_{\mathrm{prim}}\oplus \Z[E],
 \qquad \beta_h^2=2h-2,
\]
and let $\Mred(X;\gamma)$ be the reduced-virtual-class DT moduli stack in
class $\gamma$, quotienting the free elliptic translation group
$\Etr\simeq E$ and carrying the Kontsevich--Soibelman vanishing-cycle
sheaf $\phi_{W_X}$ of the CY$_3$ potential. Residual equivariance, when
used, is the chosen reduced torus/stable-automorphism equivariance after
this quotient.

Write
\[
 q=e^{2\pi i\tau},\qquad t=e^{2\pi iz},\qquad p=e^{2\pi i\sigma},
 \qquad
 \phiZO(\tau,z)=\sum_{n,\ell}c_1(4n-\ell^2)q^n t^\ell .
\]
In the K3 elliptic-genus normalisation used by the
Gritsenko--Nikulin denominator,
\[
 c_1(0)=10,\qquad
 \DFive=64^{-1}\DeltaFive,\qquad
 \PhiTenUn=\DeltaFive^2,\qquad
 \PhiTenOP=\DFive^2=4096^{-1}\DeltaFive^2 .
\]
The symbol $\PhiTenUn$ denotes the unnormalised Igusa square. The symbol
$\PhiTenOP$ denotes the monic Oberdieck--Pandharipande reduced-DT scalar
normalisation.

\begin{theorem}[Primitive $N=1$ character theorem]
\label{thm:n1-character-k3e}
For the primitive $K3$ sector of $K3\times E$, the reduced DT scalar and
the primitive BKM denominator satisfy
\[
 Z^{\mathrm{red}}_{\DT,\mathrm{prim}}(X)(q,t,p)
 =-\bigl(\PhiTenOP(\tau,z,\sigma)\bigr)^{-1}
 =-\DFive(\tau,z,\sigma)^{-2}
 =-4096\,\DeltaFive(\tau,z,\sigma)^{-2}.
\]
The theta-normalised primitive denominator is $\DeltaFive$; its monic
section is the Borcherds product
\[
 \DFive(Z)
 =e^{-2\pi i(\rho,Z)}
  \prod_{\alpha=(n,\ell,m)\in\Delta_+}
  \bigl(1-e^{-2\pi i(\alpha,Z)}\bigr)^{c_1(4nm-\ell^2)} ,
\]
after the Gritsenko--Nikulin choice of positive cone and monic leading
coefficient. Here $e^{-2\pi i(\alpha,Z)}=q^n t^\ell p^m$; in the
Gritsenko--Nikulin lattice convention this is the same denominator
written as $\DFive(2Z)$. Hence the OP/DT scalar records the doubled
character
\[
 -Z^{\mathrm{red}}_{\DT,\mathrm{prim}}(X)=\DFive^{-2}
 =
 e^{4\pi i(\rho,Z)}
 \prod_{\alpha\in\Delta_+}
 \bigl(1-e^{-2\pi i(\alpha,Z)}\bigr)^{-2c_1(4nm-\ell^2)} .
\]
The primitive Borcherds weight is
\[
 \kBKM(\DeltaFive)
 =\kBKM(\DFive)
 =\frac{c_1(0)}{2}
 =\frac{10}{2}
 =5.
\]
The square $\PhiTenUn=\DeltaFive^2$ has weight $10$; the scalar
normalisation $\PhiTenOP=4096^{-1}\DeltaFive^2$ has the same weight.
\end{theorem}

\begin{proof}
Oberdieck--Pandharipande and Oberdieck identify the reduced
$K3\times E$ primitive PT/DT scalar with the signed reciprocal of the
OP-normalised Igusa form. The OP convention uses the monic Borcherds
product $\DFive=64^{-1}\DeltaFive$, so inversion gives
\[
 -(\PhiTenOP)^{-1}=-\DFive^{-2}=-4096\,\DeltaFive^{-2}.
\]
Gritsenko--Nikulin construct $\gDeltaFive$ from the Borcherds lift of
the K3 weak Jacobi form $\phiZO$; the denominator section is the
monic product $\DFive$ with signed root multiplicities
$c_1(4nm-\ell^2)$. Borcherds' singular-theta weight formula gives
weight $c_1(0)/2$, and the K3 seed has $c_1(0)=10$. Thus the primitive
BKM denominator has weight $5$, while the OP scalar is the square of
the primitive denominator character, with the inverse-square Weyl
monomial $e^{4\pi i(\rho,Z)}$. The factor $4096=64^2$ and the minus sign
belong to the OP/DT scalar convention, leaving the primitive
BKM weight unchanged.
\end{proof}

\subsection*{Positive half and PBW primitives}

The Hall-side object at $d=3$ is $E_1$-associative. A compact oriented
critical-CoHA realisation has positive half
\[
 \Ypluscy
 =
 \HBM_{\Tred\times\AutS(K3)}
 \bigl(\Mbps(X)_{\mathrm{prim}}/\Etr,\,
       \phi_{W_X}\otimes\mathscr L^{1/2}\bigr),
\]
graded by the effective cone $\NE(X)$. When the cohomological
integrality and PBW hypotheses hold for this oriented critical CoHA, the
Davison--Meinhardt theorem supplies a PBW primitive Lie superalgebra
\[
 \operatorname{Prim}_{\mathrm{PBW}}\Ypluscy=\gBPS^+(X),
 \qquad
 \Ypluscy\simeq U(\gBPS^+(X))
\]
as an $E_1$ Hall algebra statement. If the Hall orientation chooses the
single monic square root $\DFive^{-1}$ of the OP character, then
\[
 \operatorname{sdim}\gBPS^+_{\alpha}(X)
 =
 c_1(4nm-\ell^2)
 =
 \smult_{\gDeltaFive}(\alpha),
 \qquad
 \alpha=(n,\ell,m).
\]
The OP scalar alone determines the doubled exponent
$2c_1(4nm-\ell^2)$. The primitive exponent requires this square-root
orientation and integration datum.

The compact Hall--Drinfeld upgrade requires the following data:
\begin{itemize}
\item a multiplicative orientation/Pfaffian line on the reduced
      $(-1)$-shifted symplectic moduli stack, compatible with the
      Maulik--Toda Behrend sign and the Borcherds character
      $\nu_{\DeltaFive}$;
\item a Serre-dual negative half and a nondegenerate Hopf pairing;
\item a coproduct compatible with Hall multiplication after completion;
\item a completed Borcherds-cone topology for infinite positive roots;
\item the radical quotient and PBW filtration comparison;
\item the bracket and Borcherds--Serre relation comparison with
      $\gDeltaFive$;
\item centre compatibility for the $E_1$-chiral specialisation.
\end{itemize}
These data are exactly the passage from the character theorem to an
isomorphism of completed Hall--Drinfeld/BKM objects.

\subsection*{Layer separation}

The positive Hall half, the doubled Hall object, and the chiral vertex
specialisation are separate layers:
\[
 \Yplus(X)
 \quad\rightsquigarrow\quad
 D(\Yplus(X))
 \quad\rightsquigarrow\quad
 \Phi_3(X) \text{ after the } (\Sigma_2,C)\text{-specialisation}.
\]
For $d=3$ the output algebra $A=\Phi_3(X)$ is $E_1$-chiral.
The derived centre carries the $E_2$-braided structure:
\[
 Z(\mathrm{Rep}^{E_1}(A)).
\]
The model case
$\CoHA(\C^3)=Y^+(\widehat{\mathfrak{gl}}_1)$ is a positive-half
statement. The algebra $\mathcal{W}_{1+\infty}$ appears after Drinfeld
doubling and evaluation. For $K3\times E$, the completed
Hall--Drinfeld/Borcherds branch is attached to $\DeltaFive$. The Mukai
self-mirror K3 Yangian branch is separate.

\subsection*{Subscripted invariants on $K3\times E$}

The four-value spectrum is
\[
 \{0,3,5,24\}.
\]
The zero value is carried by two different total-space invariants:
\[
 \kcat(K3\times E)
 =
 \chi(\cO_{K3})\chi(\cO_E)
 =
 2\cdot 0
 =
 0,
\]
and
\[
 \kch(K3\times E)
 =
 \sum_q(-1)^q h^{0,q}(K3\times E)
 =
 1-1+1-1
 =
 0.
\]
The Heisenberg specialisation is a different construction:
\[
 \kchHeis(K3\times E)=2+1=3.
\]
The automorphic denominator invariant is
\[
 \kBKM(\DeltaFive)=c_1(0)/2=5,
\]
and the fibre invariant is
\[
 \kfib(K3)=\operatorname{rk}\widetilde H(K3,\Z)=1+22+1=24.
\]
Adding the compact Hodge term to the elliptic fibre Euler term gives
$0+0=0$ at $N=1$, while the Borcherds theorem gives $5$. The invariant
$\kBKM$ is therefore read from the Fourier constant $c_1(0)$ of the
Jacobi input, independent of compact Hodge supertrace and fibre Euler
terms. No additive law relates the four entries.

\subsection*{The exact $N=1$ theorem}

The unconditional $N=1$ theorem consists of the OP/DT scalar identity,
the Gritsenko--Nikulin--Borcherds denominator, the root-multiplicity
formula for $\gDeltaFive$, and the weight computation
$\kBKM(\DeltaFive)=5$. The conditional Hall statement is an implication:
an oriented compact critical CoHA whose integration map chooses the
monic square root $\DFive^{-1}$ has PBW primitive superdimensions equal
to the root multiplicities of $\gDeltaFive$. The isomorphism
$\gBPS(X)\cong\gDeltaFive$ as completed Lie superalgebras is the
Hall--Drinfeld recognition problem: orientation, negative half, Hopf
pairing, coproduct, completion, radical quotient, PBW comparison,
Borcherds--Serre brackets, and centre compatibility must all be
constructed.
